% ============================================================
%  Příprava na maturitu – Mechatronika
% ============================================================

\documentclass[a4paper,10pt]{article}

% --- Balíčky ---
\usepackage[T1]{fontenc}
\usepackage[utf8]{inputenc}
\usepackage[czech]{babel}
\usepackage[left=2.4cm, right=1.8cm, top=1.8cm, bottom=1.8cm]{geometry}
\usepackage{lmodern}
\usepackage{microtype}
\usepackage{parskip}
\usepackage{titlesec}
\usepackage{enumitem}
\usepackage{xcolor}
\usepackage{hyperref}
\usepackage{tabularx}
\usepackage{array}
\usepackage{tikz}
\usepackage{eso-pic}
\usepackage{mdframed}
\usepackage{amsmath}
\usepackage{amssymb}

% --- Barvy ---
\definecolor{accent}{HTML}{03254E}
\definecolor{stripeblue}{HTML}{568EA3}
\definecolor{medgray}{HTML}{888888}
\definecolor{lightblue}{HTML}{EAF2F8}
\definecolor{lightyellow}{HTML}{FDFBE4}
\definecolor{lightgreen}{HTML}{EAF7EA}

% --- Levý pruh ---
\AddToShipoutPictureBG{
  \begin{tikzpicture}[remember picture, overlay]
    \fill[stripeblue]
      (current page.north west)
      rectangle
      ([xshift=10mm] current page.south west);
    \fill[accent!60]
      ([xshift=10mm] current page.north west)
      rectangle
      ([xshift=12mm] current page.south west);
  \end{tikzpicture}
}

\titleformat{\section}
  {\large\bfseries\color{accent}}
  {\thesection.}{0.5em}{}[\color{accent}\titlerule]

\titlespacing{\section}{0pt}{16pt}{6pt}

% --- Boxy ---
\newmdenv[backgroundcolor=lightblue,linecolor=stripeblue]{pojmybox}
\newmdenv[backgroundcolor=lightyellow,linecolor=accent!40]{vysvetlenibox}
\newmdenv[backgroundcolor=lightgreen,linecolor=accent!30]{vzorcebox}

% ============================================================
\begin{document}

\section{Nové materiály a technologie v mechatronice}

\begin{pojmybox}
\textbf{Klíčové pojmy:}
kompozity, uhlíková vlákna, SMA, piezoelektrika,
MEMS, nanomateriály, 3D tisk,
FDM, SLA, SLS, DMLS, smart materiály
\end{pojmybox}

\begin{vysvetlenibox}

% ------------------------------------------------------------
\textbf{Konstrukční materiály -- srovnání}

\begin{center}
\begin{tabular}{l r r r r}
Materiál & $\rho$ & $E$ & $R_m$ & $E/\rho$ \\
 & [g/cm³] & [GPa] & [MPa] & \\
\hline
Ocel S235 & 7,85 & 210 & 360 & 26,8 \\
Al 7075   & 2,81 & 72  & 570 & 25,6 \\
CFRP      & 1,55 & 135 & 1500 & 87,1 \\
Titan     & 4,43 & 114 & 950 & 25,7 \\
\end{tabular}
\end{center}

\textbf{Shrnutí:}
\begin{itemize}
\item Kompozity (CFRP) mají nejlepší poměr pevnost/hmotnost
\item Ocel je levná a pevná, ale těžká
\item Hliník je lehký, ale méně tuhý
\end{itemize}

% ------------------------------------------------------------
\textbf{Smart materiály}

\begin{itemize}
\item \textbf{SMA (tvarová paměť)}
    \begin{itemize}
    \item materiál si „pamatuje“ původní tvar
    \item při zahřátí se vrací zpět
    \item použití: aktuátory, medicína
    \end{itemize}

\item \textbf{Piezoelektrika}
    \begin{itemize}
    \item tlak → napětí (senzor)
    \item napětí → deformace (aktuátor)
    \item velmi rychlá odezva
    \end{itemize}

\begin{vzorcebox}
Piezo jev: $Q = d \cdot F$
\end{vzorcebox}

\item \textbf{Magnetoreologické kapaliny}
    \begin{itemize}
    \item mění viskozitu v magnetickém poli
    \item použití: tlumiče (auta, robotika)
    \end{itemize}
\end{itemize}

% ------------------------------------------------------------
\textbf{MEMS (Micro-Electro-Mechanical Systems)}

Miniaturní mechanické struktury na čipu:

\begin{itemize}
\item akcelerometry (mobil, drony)
\item gyroskopy
\item tlakové senzory
\end{itemize}

\textbf{Výhody:}
\begin{itemize}
\item malé rozměry
\item nízká cena
\item integrace s elektronikou
\end{itemize}

% ------------------------------------------------------------
\textbf{Nanomateriály}

\begin{itemize}
\item struktury v měřítku nm ($10^{-9}$ m)
\item extrémní vlastnosti:
\begin{itemize}
\item vysoká pevnost
\item vysoká vodivost
\end{itemize}
\item příklad: grafen, nanotrubice
\end{itemize}

% ------------------------------------------------------------
\textbf{3D tisk (aditivní výroba)}

\begin{center}
\begin{tabularx}{0.95\linewidth}{l X l}
Metoda & Princip & Materiál \\
\hline
FDM  & tavení plastu & PLA, ABS \\
SLA  & UV vytvrzení & pryskyřice \\
SLS  & spékání prášku & nylon \\
DMLS & tavení kovu & ocel, titan \\
\end{tabularx}
\end{center}

\textbf{Výhody:}
\begin{itemize}
\item složité tvary bez obrábění
\item rychlé prototypování
\end{itemize}

\textbf{Nevýhody:}
\begin{itemize}
\item horší mechanické vlastnosti (hlavně FDM)
\item vyšší cena u kovového tisku
\end{itemize}

% ------------------------------------------------------------
\textbf{Shrnutí}

\begin{itemize}
\item Kompozity → lehké a pevné konstrukce
\item Smart materiály → senzory + aktuátory
\item MEMS → miniaturizace
\item 3D tisk → rychlá výroba
\end{itemize}

\end{vysvetlenibox}

\end{document}
